\documentclass[11pt,a4paper]{report}
% \usepackage[none]{hyphenat} % prevents hyphenation at the end of the line
\usepackage[left=1.9in,top=1.5in,right=1.6in,bottom=1.5in]{geometry} % geometry package is used to easily set all the margins
\usepackage{graphicx}
\usepackage{tabularx}
\usepackage{ragged2e}   % this is used to justify any type of text
\usepackage{tikz}   % package used to draw page borders
\newcommand{\HRule}{\rule{\linewidth}{0.6mm}}	% this sets the thickness of a horizontal line



%=========================must use in every document u prepare=====================
\linespread{1.3}    % 1.5 line spacing
\addtolength{\oddsidemargin}{-0.4in}
\addtolength{\evensidemargin}{0.4in}
\addtolength{\textwidth}{1in}
\addtolength{\topmargin}{-2.5cm}
\addtolength{\textheight}{1.4in}
%=========================must use in every document u prepare=====================


\begin{document}

%===================================title page===================================
\begin{titlepage}

\begin{center}

\fontsize{12}{12}	% font size of the paragraph text
\selectfont

% page border using tikz package
\begin{tikzpicture}[remember picture,overlay]
% drawing the outer line
    \draw[line width=0.5mm] ([xshift=2.7cm,yshift=-1.3cm]current page.north west) rectangle ([xshift=-1.6cm,yshift=1.3cm]current page.south east);
% now drawing the inner line
    \draw[line width=0.2mm] ([xshift=3cm,yshift=-1.6cm]current page.north west) rectangle ([xshift=-1.9cm,yshift=1.6cm]current page.south east);
\end{tikzpicture}



J S S Mahavidyapeetha \\
\textbf{Sri Jayachamarajendra College of Engineering (SJCE)\\ Mysore - 570 006} \\
An Autonomous Institute Affiliated to \\
Visvesvaraya Technological University, Belgaum \\
\end{center}
\smallskip



\hspace{0.3cm}
\includegraphics[width=0.14\textwidth]{./VTU_logo.png} 
\hspace{8cm}
\includegraphics[width=0.20\textwidth]{./sjce.jpg} 	\\[0.5cm]



\begin{center}
\textbf{CERTIFICATE} \\[1cm]
\end{center}



{\justifying
\hspace{0.9in}
This is to certify that the work entitled \textbf{``Support Vector Machine Approach To Automated Cardiac SPECT Diagnosis"} is a bonafied work carried out by \textbf{Dharmatheja Bhat, Supreeth Nag V.P, Varun B Patil, Vinay M} in partial fulfillment of the award of the degree of \textbf{Bachelor of Engineering in Computer Science and Engineering of Visvesvaraya Technological University, Belgaum during the year 2012.} It is certified that all corrections / suggestions indicated during CIE have been incorporated in the report. The project report has been approved as it satisfies the academic requirements in respect of project work prescribed for the Bachelor of Engineering degree.}\\
\vfill


\hspace{-0.9cm}	% pull the whole table to the left a little	
% with tabular* , u can set the width of the table - in this case = \textwidth
\begin{tabular*}{\textwidth}{lr}
% with \makebox, you cannot introduce a linebreak within a cell, but with \parbox, u can specify \newline
% but \parbox does not have left and right justification like in \makebox... instead we can use \flushleft 
% and \flushright 
\parbox{9.5cm}{\emph{\textbf{Guide}} \newline \textbf{Dr. M.P Pushpalatha} \newline Associate Professor, \newline Dept. of CS \& E, \newline S.J.C.E, Mysore} 
& \parbox{5cm}{\emph{\textbf{Head of the Department}} \newline \textbf{Dr. C.N. Ravikumar} \newline Professor \& Head, \newline Dept. of CS \& E, \newline S.J.C.E, Mysore} \\
\end{tabular*}
\vspace{1.5cm}



\hspace{-0.9cm}	% pull the whole table to the left a little	
% with tabular* , u can set the width of the table - in this case = \textwidth
\begin{tabular*}{\textwidth}{lcr}
% with \makebox, you cannot introduce a linebreak within a cell, but with \parbox, u can specify \newline
% but \parbox does not have left and right justification like in \makebox... instead we can use \flushleft 
% and \flushright 
\parbox[t]{7cm}{\flushleft \textbf{Place:} Mysore \newline \textbf{Date:}} &
\parbox[t]{3cm}{\flushright \textbf{Examiners:}} &
\parbox[t]{7cm}{\flushleft 1. .............................. \newline \newline 2. .............................. \newline \newline 3. .............................. \newline } \\
\end{tabular*}


\end{titlepage}
%==================================end of title page=========================

\end{document}
